%%%%%%%%%%%%%%%%%%%%%%%%%%%%%%%%%%%%%%%%%%%%%%%%%%%%%%%%%%%%%%%%%%%%%%%%%%%%%%%%

\pagestyle{empty}
\begin{abstract}
  Este Trabajo de Fin de Grado presenta el diseño y desarrollo de un sistema de Business Intelligence orientado a apoyar la toma de decisiones de la administración educativa autonómica en materia de distribución de recursos para la prevención del abandono escolar temprano. El sistema integra datos académicos, demográficos y socioeconómicos procedentes de sistemas transaccionales, y los transforma mediante una capa ETL implementada en Apache Hop hacia un Datamart en PostgreSQL con esquema en estrella. Durante este proceso se calculan índices de riesgo de abandono que combinan factores académicos (suspensos, repeticiones, desfase de edad), de adaptación curricular y socioeconómicos.

  La información resultante se expone a través de una API REST en FastAPI y se visualiza en un \textit{dashboard} analítico en React con dos niveles de análisis: una vista Macro que permite a los técnicos de la administración identificar y comparar centros educativos a escala autonómica, y una vista Micro que posibilita el diagnóstico causal del riesgo de un centro concreto, facilitando decisiones de asignación de recursos fundamentadas en evidencia. Todo el ecosistema se despliega mediante contenedores Docker, garantizando su portabilidad y reproducibilidad.

  \vspace*{25pt}
  \begin{segundoresumo}
    This Bachelor's Thesis presents the design and development of a Business Intelligence system aimed at supporting data-driven decision-making by regional educational authorities in the allocation of resources to prevent early school dropout. The system integrates academic, demographic, and socioeconomic data from transactional sources and transforms them through an ETL layer built with Apache Hop into a star-schema Datamart on PostgreSQL. During this process, dropout risk indices are computed by combining academic factors (failures, grade repetition, age gap), curricular adaptation needs, and socioeconomic indicators.

    The resulting information is exposed through a FastAPI REST API and visualised in a React analytical dashboard with two levels of analysis: a Macro view enabling regional administrators to identify and compare schools at an autonomous-community scale, and a Micro view allowing causal diagnosis of the risk profile of a specific school, supporting evidence-based resource allocation decisions. The entire ecosystem is deployed using Docker containers, ensuring portability and reproducibility.
  \end{segundoresumo}
\vspace*{25pt}
\begin{multicols}{2}
\begin{description}
\item [\palabraschaveprincipal:] \mbox{} \\[-20pt]
  \begin{itemize}
    \item Inteligencia de Negocios
    \item Prevención Abandono Escolar
    \item Datamart
    \item ETL
    \item Apache Hop
    \item Contenedores Docker
  \end{itemize}
\end{description}
\begin{description}
\item [\palabraschavesecundaria:] \mbox{} \\[-20pt]
  \begin{itemize}
    \item Business Intelligence
    \item School Dropout Prevention
    \item Datamart
    \item ETL
    \item Apache Hop
    \item Docker Containers
  \end{itemize}
\end{description}
\end{multicols}

\end{abstract}
\pagestyle{fancy}

%%%%%%%%%%%%%%%%%%%%%%%%%%%%%%%%%%%%%%%%%%%%%%%%%%%%%%%%%%%%%%%%%%%%%%%%%%%%%%%%
